% chapters/ch20-constraints-catalog.tex
\chapter{Constraints Catalog}

\begin{learningobjectives}
In this chapter you will learn a structured catalog of constraints.
You will distinguish uniqueness, inclusion, and domain constraints.
You will practice expressing constraints without SQL.
\end{learningobjectives}

\section{Why a catalog helps}
Constraints define valid states.
A catalog ensures you do not miss common rule types.
It also improves review discussions.

\section{Uniqueness constraints}
Uniqueness defines identity and prevents duplicates.
It can apply to a single attribute or a combination.

\subsection{Single-attribute uniqueness}
Example: a member is uniquely identified by \texttt{member\_id}.
This is a primary identity rule.

\subsection{Composite uniqueness}
Composite uniqueness captures meaning like “no double booking.”
It is a relationship-level rule expressed as a unique combination.

\begin{examplebox}
\textbf{Composite uniqueness}\\
Rule: A doctor cannot have two appointments at the same start time.\\
Meaning: (doctor\_id, start\_time) is unique in \rel{Appointment}.
\end{examplebox}

\section{Inclusion dependencies}
Inclusion dependencies express referential integrity.
They say one set of values must be contained in another.

\subsection{Mandatory references}
If a loan references a member, the member must exist.
This is a core relational meaning, not a storage detail.

\begin{examplebox}
\textbf{Inclusion dependency}\par\smallskip
\[
\rel{Loan}.member\_id \rightarrow \rel{Member}.member\_id
\]
The referenced member must exist.
\end{examplebox}

\section{Domain constraints}
Domain constraints restrict allowed values.
They protect meaning at the attribute level.

\subsection{Enumerations}
Status values should come from a defined set.
This prevents silent drift in meaning.

\subsection{Value ranges}
Dates, counts, and measures often have valid ranges.
These are cheap but powerful constraints.

\begin{principlebox}
\textbf{Principle}\\
If a value has meaning, it has constraints.
\end{principlebox}

\section{Cross-row constraints}
Some rules depend on multiple rows.
They are often business constraints.
They must still be stated conceptually.

\subsection{No overlap}
Examples include no overlapping reservations or appointments.
State the rule in words even if enforcement is complex.

\section{Temporal constraints}
Time introduces ordering and validity windows.
Rules like “return date after checkout date” are temporal.
They should be explicit in the model.

\begin{chaptersummary}
Constraints are not optional.
Use a catalog to ensure you cover uniqueness, inclusion, and domain rules.
Cross-row and temporal constraints preserve real-world meaning.
\end{chaptersummary}

\begin{exercisebox}
Build a constraint list for a small domain.
\begin{enumerate}[label=\arabic*.]
  \item Write two uniqueness constraints (one composite).
  \item Write two inclusion dependencies.
  \item Write two domain constraints (enum or range).
  \item Write one cross-row constraint.
  \item Explain how each constraint protects meaning.
\end{enumerate}
\end{exercisebox}
